\section{Transfer to finite factors}
\label{sec:finite-factors}

We now transfer Theorem~\ref{thm:free-corner} to matrix commutators in
finite sums of the standard representation of a separable-predual
\(\mathrm{II}_1\) factor.  All traces on matrix amplifications are normalized.

We first average the operator over an irreducible subfactor generated by
an increasing sequence of full matrix algebras. After matrix amplification,
the ultrapower of the amplified subfactor contains another subfactor
generated by increasing full matrix algebras and free from the constant
copy of the amplified ambient factor. Extending and compressing the
averaged operator then puts it in the setting of
Theorem~\ref{thm:free-corner}. Compression to constant vectors transfers
the estimate back to the matrix test, with the fixed error from the
initial averaging.

The existence of an irreducible hyperfinite subfactor in a
separable-predual \(\mathrm{II}_1\) factor is due to Popa
\citep{Popa1981}; we include a direct proof of the increasing-matrix
form used below.

\begin{lemma}[An irreducible increasing-matrix subfactor]
\label{lem:irreducible-hyperfinite}
Every separable-predual \(\mathrm{II}_1\) factor \(M\) contains a unital
subfactor \(R\) with \(R'\cap M=\C1\), generated by an increasing sequence
of unital full matrix algebras whose sizes tend to infinity.
\end{lemma}

\begin{proof}
We first give an approximation step in an arbitrary \(\mathrm{II}_1\)
factor \(N\).  If \(y=y^*\in N\) and \(\eta>0\), there is a unital full
matrix algebra \(F\subset N\), of size at least two, such that
\[
 \norm{E_{F'\cap N}(y)-\tau(y)1}_2<\eta.
\]
Indeed, approximate \(y\) in \(L^2\) by a spectral step
\(y_0=\sum_{i=1}^k\lambda_i p_i\), where the \(p_i\) partition the
identity.  For a large integer \(b\), split off
\(\lfloor b\tau(p_i)\rfloor\) subprojections of trace \(1/b\) in each
\(p_i\).  The remaining projection has trace
\[
 \frac{b-\sum_{i=1}^k\lfloor b\tau(p_i)\rfloor}{b}<\frac{k}{b}.
\]
Its numerator is an integer, so it too splits into projections of trace
\(1/b\).  These pieces form a partition of the identity by \(b\)
equivalent projections and hence extend to matrix units for a unital
\(F\cong M_b(\C)\).  Define \(y_1\in F\) to have value \(\lambda_i\)
on the pieces taken from \(p_i\), and value zero on the remaining pieces.
Then
\[
 \norm{y-y_1}_2
 \leq \norm{y-y_0}_2+\norm{y_0}\sqrt{k/b}.
\]
The expectation of \(y_1\) onto \(F'\cap N\) is \(\tau(y_1)1\).
Contractivity of this expectation on \(L^2\), and
\(\abs{\tau(y-y_1)}\leq\norm{y-y_1}_2\), give the required estimate
by choosing the spectral approximation and then \(b\).

Choose a countable \(L^2\)-dense subset \(\{x_j:j\geq1\}\) of the
self-adjoint unit ball of \(M\), and enumerate all requirements
\((j,k)\in\N^2\).  Begin with \(F_0=\C1\).  If the current stage is
a unital full matrix algebra \(F_n\), then
\(N_n=F_n'\cap M\) is a \(\mathrm{II}_1\) factor: matrix units identify
\(M\) with a matrix algebra over a corner and \(N_n\) with that corner.
For the next requirement \((j,k)\), apply the approximation step in
\(N_n\) to \(y=E_{N_n}(x_j)\), with error less than \(1/k\), obtaining
\(F\subset N_n\).  Set \(F_{n+1}=F_n\vee F\).  This is again a full
matrix algebra, its size is at least twice that of \(F_n\), and
\(N_{n+1}=F'\cap N_n\).  The identity
\(E_{N_{n+1}}=E_{N_{n+1}}E_{N_n}\) gives
\[
 \norm{E_{N_{n+1}}(x_j)-\tau(x_j)1}_2<1/k.
\]
All previously obtained bounds persist because the relative commutants
decrease and their trace expectations compose.

Let \(R=(\bigcup_nF_n)''\) and \(N=R'\cap M=\bigcap_nN_n\).
Since \(E_N=E_NE_{N_n}\), every requirement remains valid with \(N\)
in place of \(N_{n+1}\).  Thus \(E_N(x_j)=\tau(x_j)1\) for every
\(j\).  Density gives \(E_N(x)=\tau(x)1\) for all \(x\in M\), so
\(N=\C1\).  In particular \(R\) is a factor.  It is finite and contains
full matrix algebras of unbounded size, hence is of type
\(\mathrm{II}_1\), with the asserted increasing-matrix presentation.
\end{proof}

For a finite tracial von Neumann algebra $P$ and a nonprincipal
ultrafilter $\omega$ on $\N$, write
$P^\omega=\ell^\infty(\N,P)/\mathcal I_\omega$, where
$\mathcal I_\omega$ consists of the uniformly bounded sequences
$(a_i)$ with $\lim_\omega\|a_i\|_2=0$. Its trace is
$\tau_\omega([a_i])=\lim_\omega\tau(a_i)$, and $P$ embeds by
constant sequences. The lemma below places a free subfactor in such
an ultrapower while retaining its inclusion in the prescribed
subalgebra $Q^\omega$.

\begin{lemma}[Free position in the full ultrapower]
\label{lem:free-position}
Let \(Q\subset P\) be a unital irreducible inclusion of
\(\mathrm{II}_1\) factors, with \(P\) of separable predual.  For any
nonprincipal ultrafilter \(\omega\) on \(\N\), there is a subfactor
\(D\subset Q^\omega\), generated by increasing unital full matrix
algebras of unbounded sizes, which is free from the constant copy of
\(P\) in the tracial ultrapower \(P^\omega\).
\end{lemma}

\begin{proof}
We use the ultraproduct-of-subalgebras case of Popa's independence
theorem \citep[Theorem~0.1(a); see also Theorem~4.3]{Popa2014}.
For von Neumann subalgebras $A,B\subset M$, the notation
$A\prec_M B$ means Popa's intertwining of a nonzero corner of $A$
into a finite matrix amplification of $B$.
Its hypotheses are finite factors \(M_i\) with
\(\dim M_i\longrightarrow\infty\), and unital von Neumann subalgebras
\(Q_i\subset M_i\) satisfying
\(Q_i\nprec_{M_i}(Q_i'\cap M_i)\).  Write
\(\mathbf M=\prod_\omega M_i\), \(\mathbf Q=\prod_\omega Q_i\), and
\(C=\mathbf Q'\cap\mathbf M\).  Given an \(L^2\)-separable subspace
\(X\subset\mathbf M\ominus C\), the theorem supplies a diffuse abelian
von Neumann algebra \(A\subset\mathbf Q\) such that
\[
 E_C(x_0a_1x_1\cdots a_nx_n)=0
\]
whenever \(n\geq1\), \(x_0\in X\cup\{1\}\), \(x_i\in X\) for
\(1\leq i\leq n\), and \(a_i\in A\ominus\C1\).
In particular, if \(C=\C1\) and \(T\subset\mathbf M\) is a
separable-predual von Neumann subalgebra, take
\(X=T\ominus\C1\).  The moment identities, traciality, and decomposition
into scalar and centered parts show that \(A\) is free from \(T\).
Here separability is for \(L^2\), as stipulated in
\citet[Section~1.1]{Popa2014}.

We check these hypotheses for the constant sequences \(M_i=P\) and
\(Q_i=Q\).  The dimension condition holds because these factors are
infinite-dimensional. If \(Q\prec_P\C1\), the intertwining data
would include a nonzero normal \(*\)-homomorphism from a nonzero
corner of \(Q\) into a finite-dimensional matrix algebra.  Such a
homomorphism from a factor is faithful, whereas a
\(\mathrm{II}_1\) factor has arbitrarily large families of nonzero
orthogonal projections.  This is impossible.  Since \(Q'\cap P=\C1\),
the non-intertwining hypothesis follows.

We also verify the full ultrapower relative commutant:
\[
                 (Q^\omega)'\cap P^\omega=\C1.
\]
For a uniformly bounded representative \((x_i)\) of an element of this
commutant, commuting with all sequences of unitaries from \(Q\) implies
\[
 d_i:=\sup_{v\in\mathcal U(Q)}
             \norm{vx_iv^*-x_i}_2\longrightarrow_\omega0.
\]
Otherwise one could choose a witnessing unitary at each coordinate
on an \(\omega\)-large set.  The \(L^2\)-closed convex hull of the
\(Q\)-unitary orbit of \(x_i\) has a unique vector of minimum norm.
It is \(Q\)-central, has trace \(\tau(x_i)\), and belongs to \(P\):
the hull stays in the operator-norm ball of radius \(\norm{x_i}\),
which is \(L^2\)-closed by weak-* compactness and trace pairings.
Irreducibility makes this vector \(\tau(x_i)1\).  Every point in the
hull is within \(d_i\) of \(x_i\), so
\(\norm{x_i-\tau(x_i)1}_2\leq d_i\).  Bounded scalar sequences equal
their scalar ultralimit in the tracial ultrapower.  This proves the
claim, which concerns all of \(Q^\omega\), rather than merely the
constant copy of \(Q\).

Choose a Haar unitary \(u_1\) in the diffuse abelian algebra obtained
by applying Popa's theorem to the constant \(P\).  Then
\(u_1\in Q^\omega\) is free from \(P\).
The algebra \(W^*(P,u_1)\) has separable \(L^2\) space: countably
many words in generators of \(P\), in \(u_1\), and in their adjoints
have linear span dense in \(L^2\).
A second application of the same theorem therefore gives a Haar
unitary \(u_2\in Q^\omega\) free from \(W^*(P,u_1)\).
Expanding centered reduced words shows that
\(P,W^*(u_1),W^*(u_2)\) are jointly free.  Thus
\(W^*(u_1,u_2)\) is free from \(P\) and, by its free Haar moments,
is trace-isomorphic to \(L(\mathbb F_2)\).  This is a separable-predual
\(\mathrm{II}_1\) factor: the nontrivial conjugacy classes of
\(\mathbb F_2\) are infinite, so a central square-summable Fourier
coefficient function is supported at the identity.  Apply
Lemma~\ref{lem:irreducible-hyperfinite} inside \(W^*(u_1,u_2)\) to
obtain \(D\).  Freeness passes to its subalgebra \(D\), proving the
lemma.  The ambient tracial ultrapower throughout is \(P^\omega\).
\end{proof}

\begin{proposition}[Finite standard multiplicities]
\label{prop:finite-copies}
Let \(M\) be a separable-predual \(\mathrm{II}_1\) factor, and let
\(M\) act on \(H_r=L^2(M)^r\) by common left multiplication, where
\(r\geq1\) is finite.  For every \(Y\in\mathcal B(H_r)\), every
finite \(h\geq1\), and every \(X\in M_h(M)\),
\[
 \norm{[Y^{(h)},X]}
       \leq (3c_1+2)\,g_M(Y)\,\norm{X}.
\]
Here \(X\) denotes its entrywise action on \(H_r^h\); the constant
is independent of \(M,r,h\).
\end{proposition}

\begin{proof}
Put \(\delta=g_M(Y)\).  By homogeneity it suffices to consider
\(\norm{X}\leq1\).  We first arrange a symmetry test on a square
array of standard copies.  Choose \(h_0\geq h\) so that
\(m=2h_0\geq r\), extend \(Y\) by zero to \(H_m\), and retain the
notation \(Y\).  The copy projection reduces the common action of
\(M\), so this extension still has scalar commutator bound \(\delta\).
For \(X'=X\oplus0\in M_{h_0}(M)\), define in \(P=M_m(M)\)
\[
 s=
 \begin{pmatrix}
  (1-X'X'^*)^{1/2}&X'\\
  X'^*&-(1-X'^*X')^{1/2}
 \end{pmatrix}.
\]
Functional calculus gives
\((1-X'X'^*)^{1/2}X'=X'(1-X'^*X')^{1/2}\), and hence
\(s=s^*=s^{-1}\).  Traciality gives equal moments for
\(X'X'^*\) and \(X'^*X'\), including the zeroth moment; polynomial
approximation on \([0,1]\) gives equal traces for their displayed
square roots.  Consequently \(\tau_P(s)=0\).

Write vectors of \(H_m^m\) as arrays \((\xi_{ij})\), with \(i\)
the outer amplification coordinate and \(j\) the copy coordinate.
The map \((\xi_{ij})\mapsto\sqrt m\,(\xi_{ij})\) is a unitary from
\(H_m^m\) to \(L^2(P)\), since
\(\norm{(\xi_{ij})}_2^2=m^{-1}\sum_{i,j}\norm{\xi_{ij}}_2^2\)
on the latter space.  Under this identification, the matrix action
of \(P\) is left multiplication.  The upper-right block of
\([Y^{(m)},L_s]\) is the commutator with \(X'\).  Restricting its
row coordinates to the original \(h\) and its copy coordinates to
the original \(r\) recovers \([Y^{(h)},X]\) for the original operator.
It is therefore enough to bound \([Y^{(m)},L_s]\).

Choose \(R\subset M\) from Lemma~\ref{lem:irreducible-hyperfinite},
with matrix stages \(F_n\).  On \(H_m\), average \(Y\) over the compact
unitary group of each \(F_n\), using common left multiplication,
and take a weak operator cluster point \(Y_0\) as \(n\) tends to
infinity.  For every such unitary \(v\),
\(\norm{L_vYL_v^*-Y}\leq\delta\).  The averages, and hence their
cluster point, are within \(\delta\) of \(Y\).  An average over a
later stage commutes with every earlier stage, so
\[
 [Y_0,L_R]=0,\qquad
 \norm{Y_0-Y}\leq\delta,\qquad g_M(Y_0)\leq3\delta.
\]
The first identity extends from the matrix stages to \(R\) by strong
density; the norm bounds pass to the cluster point by weak operator
closedness of norm balls.

On \(L^2(P)=H_m^m\), set \(S=Y_0^{(m)}\) and \(Q=M_m(R)\).
Matrix units and \(R'\cap M=\C1\) imply \(Q'\cap P=\C1\).
Moreover \(S\) commutes with \(L_Q\), since it has the same operator
\(Y_0\) on every row and \(Y_0\) commutes with every entry from \(R\).
For \(e_0=e_{11}\otimes1\), the space \(e_0L^2(P)\) is one row of
\(m\) standard copies, with \(e_0Pe_0\cong M\) acting by common
left multiplication.  Since \(S\) commutes with \(L_{e_0}\), it follows
on the full Hilbert space that
\[
 \norm{[S,L_z]}\leq3\delta\norm{z}\qquad(z\in e_0Pe_0).
\]
If \(f\in Q\) is a projection with \(\tau_P(f)\leq1/m\), projection
comparison in the factor \(Q\) gives a unitary \(v\in Q\) with
\(f\leq ve_0v^*\).  Conjugating by \(L_v\), which commutes with
\(S\), yields
\[
 \norm{[S,L_z]}\leq3\delta\norm{z}
       \qquad(z\in fPf,\ f\in Q,\ \tau_P(f)\leq1/m).
\]

We have obtained the same small-corner estimate for every projection
of $Q$ of trace at most $1/m$. The next step transports that estimate
to the ultrapower, where Lemma~\ref{lem:free-position} supplies the
free subfactor. Fix a nonprincipal ultrafilter $\omega$ on $\N$.
Recall the tracial quotient and distinguish it from the Hilbert
ultrapower:
\[
 \begin{aligned}
 P^\omega&=\ell^\infty(\N,P)/\mathcal I_\omega,
 &\mathcal I_\omega&=
 \{(a_i):\lim_\omega\norm{a_i}_2=0\},\\
 \mathcal H&=(L^2(P))^\omega_{\mathrm{Hilb}}.
 \end{aligned}
\]
The Hilbert ultrapower uses bounded Hilbert-space sequences modulo
those whose Hilbert norms tend to zero along \(\omega\).  Sending
the class of a uniformly operator-bounded sequence \((a_i)\) to its
Hilbert-ultrapower class is a trace-preserving isometry on the
dense subspace \(P^\omega\subset L^2(P^\omega)\).  Its extension
identifies \(\mathcal C=L^2(P^\omega)\) with a closed subspace of
\(\mathcal H\).

For a bounded sequence \(a=(a_i)\) in \(P\), its coordinatewise left
action on \(\mathcal H\), and the adjoint action corresponding to
\((a_i^*)\), preserve the bounded-algebra-sequence subspace of
\(\mathcal C\).  Thus \(\mathcal C\) reduces both actions.  If
\(a\in\mathcal I_\omega\), then for every operator-bounded sequence
\((b_i)\),
\[
 \norm{a_ib_i}_2\leq\norm{b_i}\norm{a_i}_2
                       \longrightarrow_\omega0.
\]
It follows by density that the action on \(\mathcal C\) factors
through the tracial quotient and is precisely left multiplication
by \([a_i]\).  This factoring assertion is needed only on
\(\mathcal C\), not on the whole Hilbert ultrapower.

Let \(S^{\omega}_{\mathrm{Hilb}}\) be the coordinatewise action of
the fixed bounded operator \(S\) on \(\mathcal H\), and set
\[
 \widehat S=P_{\mathcal C}S^{\omega}_{\mathrm{Hilb}}
                              \big|_{\mathcal C}.
\]
Since \(S\) commutes with every left multiplier from \(Q\), its
Hilbert ultrapower commutes with every bounded left \(Q\)-sequence.
Compression to the reducing subspace \(\mathcal C\) proves
\([\widehat S,L_{Q^\omega}]=0\).

Set \(t=1/(2m)\).  We claim that
\begin{equation}
 \label{eq:ultrapower-corner}
 \begin{gathered}
 \norm{[\widehat S,L_z]}\leq3\delta\norm{z},\\
 e\in Q^\omega\text{ a projection},\qquad
       \tau(e)=t,\qquad z\in eP^\omega e.
 \end{gathered}
\end{equation}
To lift \(e\), take a bounded self-adjoint representative \((a_i)\)
in \(Q\) and set \(f_i=1_{[1/2,\infty)}(a_i)\).  The scalar inequality
\(\abs{x-1_{[1/2,\infty)}(x)}\leq2\abs{x^2-x}\) gives
\(\norm{f_i-a_i}_2\longrightarrow_\omega0\), so \([f_i]=e\) and
\(\tau(f_i)\longrightarrow_\omega t\).  The strict gap
\(t<1/m\) allows us to replace \(f_i\) by zero off an
\(\omega\)-large set and assume \(\tau(f_i)\leq1/m\) for every \(i\).

Given \(z\in eP^\omega e\) and \(\eta>0\), the \(C^*\)-quotient
norm gives a representative \((b_i)\) with
\(\sup_i\norm{b_i}\leq\norm{z}+\eta\).  The sequence
\(z_i=f_ib_if_i\) represents \(z\), lies in the required corners,
and has the same norm bound.  Consequently
\[
 \norm{[S,L_{z_i}]}\leq3\delta(\norm{z}+\eta)
 \qquad\text{for every }i.
\]
This estimate passes to the Hilbert ultrapower.  Its compression to
\(\mathcal C\) is \([\widehat S,L_z]\), since \(\mathcal C\) reduces
the bounded left sequence action.  Letting \(\eta\) decrease to zero
proves \eqref{eq:ultrapower-corner}.

By Lemma~\ref{lem:free-position}, choose \(D\subset Q^\omega\)
free from the constant \(P\), and put
\(\mathcal L=W^*(D,P)\subset P^\omega\).  Faithfulness of the trace
and the reduced-word inner products identify this algebra with the
tracial free product \(D*P\).  Let
\(G=P_{L^2(\mathcal L)}\widehat S|_{L^2(\mathcal L)}\).
The subspace \(L^2(\mathcal L)\) reduces left multiplication by
\(\mathcal L\).  Thus \(G\) commutes with \(L_D\), and
\eqref{eq:ultrapower-corner} gives the corner hypotheses of
Theorem~\ref{thm:free-corner} with
\(A_0=P\), \(\epsilon=3\delta\), and \(t=1/(2m)\).

Finally, every constant \(L^2(P)\)-vector belongs to \(\mathcal C\),
by approximation by bounded vectors from \(P\), and
\(S^{\omega}_{\mathrm{Hilb}}\) sends the constant vector \(\xi\) to
the constant vector \(S\xi\).  Both of these vectors also belong to
\(L^2(\mathcal L)\).  Therefore both compressions preserve the action
of \(S\) on the constant subspace.  As \(L_s\) preserves that subspace,
for \(\xi,\eta\in L^2(P)\) we have exactly
\[
 \langle[G,L_s]\xi,\eta\rangle
        =\langle[S,L_s]\xi,\eta\rangle.
\]
Theorem~\ref{thm:free-corner}, applied to the trace-zero symmetry
\(s\in P\), consequently yields on the original \(L^2(P)\)
\[
                   \norm{[S,L_s]}\leq3c_1\delta.
\]
Replacing \(Y_0\) by \(Y\) costs at most
\(2\norm{Y_0-Y}\norm{s}\leq2\delta\), since Hilbert amplification
preserves the difference norm.  The initial block compression now
gives the claimed bound \((3c_1+2)\delta\) for the original finite
\(r,h\).
\end{proof}
